\chapter{Grondslagen: hoe er gerekend wordt}
\label{h:grondslagen}

\section{Doel en leeswijzer}
Dit verslag rekent het procesflowschema (PFD) van 8 ITM Neptune V-elektrolyse-units stap voor stap door. Bij elke
stap staat \emph{wat} er fysiek gebeurt, \emph{welke} Cengel-methode erbij hoort, \emph{hoe} de getallen worden ingevuld
en \emph{waarom} de gekozen waarden zo zijn. Het doel is dat je elke vraag in de trant van ``waarom dit?'',
``hoezo deze temperatuur?'' of ``hoe werkt dit apparaat?'' kunt beantwoorden. In bijlage~\ref{b:vragen} staan
daarvoor honderden verdedigingsvragen met antwoorden.

Er bestaan twee versies van dit verslag met dezelfde opbouw:
\begin{itemize}
  \item \textbf{Versie A (CoolProp met Cengel-check):} de stofeigenschappen komen uit CoolProp. Dat is precies
        wat de visuele PFD-app doet. Elke CoolProp-opzoeking wordt uitgelegd als een tabelopzoeking à la Cengel.
  \item \textbf{Versie B (Cengel-handberekening):} alles met de hand: stoomtabellen (A-4, A-5, A-6), ideaal gas
        met constante $c_p$ (tabel A-2) en, voor de warmtepompen, het aflezen van enthalpieën uit het
        log $p$–$h$-diagram. Dat is dezelfde aanpak als Cengel hoofdstuk~11 met het $P$–$h$-diagram van R-134a.
\end{itemize}
%% if A
\textbf{Je leest nu versie A.} Waar CoolProp en de handmethode verschillen, staat het verschil met de oorzaak
erbij in een paars kader ``CoolProp-aanroep en Cengel-check''.
%% else
\textbf{Je leest nu versie B.} De diagrammen van de koudemiddelen zijn getekend met CoolProp, omdat Cengel geen
tabellen heeft voor R1234ze(E) en R1233zd(E). De getallen lees je daarna uit dat diagram af, net zoals je dat
in Cengel met figuur A-14 (R-134a) zou doen.
%% endif

\textbf{Bedrijfspunt.} Alles is doorgerekend voor de basisinvoer van het model (bijlage~\ref{b:invoer}):
\VAR{I.n_units} units op 100\,\% last en EOL-warmteproductie (end of life). Het warmtenet en
warmtepomp WP-201 staan aan, de industrie vraagt \VAR{n(I.Q_ind,0)}\kW. De stromen in de \OO-lijn,
de \HH-lijn en de waterkant gelden \textbf{per unit}. De warmteketen en de warmtepompen gelden voor de
\textbf{hele installatie} (8 units samen). Waar met 8 vermenigvuldigd wordt, staat dat er altijd expliciet bij.

\section{Eenheden en omrekeningen}
In dit verslag gebruiken we:
\begin{center}
\begin{tabular}{p{5.4cm}p{3.2cm}p{6cm}}
\toprule
grootheid & eenheid & opmerking\\
\midrule
massadebiet $\md$ & kg/h & \emph{altijd}; nooit kg/s of t/h\\
vermogen, warmtestroom $\Qd,\Wd$ & kW $=$ kJ/s & \\
specifieke enthalpie $h$ & kJ/kg & \\
specifieke entropie $s$, warmtecapaciteit $c_p$ & kJ/(kg$\cdot$K) & \\
druk $P$ & kPa absoluut & barg tussen haakjes, zoals op de PFD\\
temperatuur & \textdegree C, in formules met $T$ in K & $T[\mathrm{K}] = T[^{\circ}\mathrm{C}] + 273{,}15$\\
volume, volumestroom & m$^3$, m$^3$/h & \\
\bottomrule
\end{tabular}
\end{center}

\kop{De omrekening kg/h $\rightarrow$ kW}
Cengel schrijft de energiebalans met $\md$ in kg/s. Omdat wij in kg/h werken, staat er in elke energieberekening
een deling door 3600\,s/h:
\[
\Qd\,[\mathrm{kW}] \;=\; \frac{\md\,[\mathrm{kg/h}]}{3600\,[\mathrm{s/h}]}\cdot \Delta h\,[\mathrm{kJ/kg}]
\qquad\Longleftrightarrow\qquad
\md\,[\mathrm{kg/h}] = \frac{\Qd\,[\mathrm{kW}]\cdot 3600}{\Delta h\,[\mathrm{kJ/kg}]}.
\]
\textbf{Controle van de eenheden:} $\frac{\mathrm{kg/h}}{\mathrm{s/h}}\cdot\frac{\mathrm{kJ}}{\mathrm{kg}} = \frac{\mathrm{kJ}}{\mathrm{s}} = \mathrm{kW}$.

\kop{De omrekening barg $\rightarrow$ kPa absoluut}
Een manometer meet overdruk ten opzichte van de atmosfeer. Thermodynamische tabellen en de ideale-gaswet werken
met \emph{absolute} druk (Cengel H1, ``gage and absolute pressure''). Met 1\,bar $=$ 100\,kPa en
$P_\mathrm{atm} = 101{,}325$\,kPa:
\[
P_\mathrm{abs}\,[\mathrm{kPa}] = P_\mathrm{g}\,[\mathrm{barg}]\cdot 100 + 101{,}325 .
\]
Voorbeeld: de \HH-compressor perst naar 70\,barg, dus $P = 70\cdot100+101{,}325 = \VAR{n(R.P_uit_K101,1)}\kPa$.

\section{Het toestandspostulaat: twee eigenschappen leggen alles vast}
\begin{methode}[– Cengel H3, ``The state postulate'']
De toestand van een \emph{enkelvoudig samendrukbaar systeem} (een zuivere stof zonder elektrische, magnetische
of oppervlakte-effecten) ligt vast met \textbf{twee onafhankelijke intensieve eigenschappen}. Ken je er twee,
dan liggen alle andere vast: $h$, $s$, $v$, $u$, enzovoort.
\end{methode}
Dit is de basis van \emph{elke} opzoeking in dit verslag, met de hand of met CoolProp:
\begin{itemize}
  \item \textbf{Eenfasegebied} (oververhitte damp, gas, onderkoelde vloeistof): $T$ en $P$ zijn onafhankelijk.
        Je gebruikt bijvoorbeeld tabel A-6 met $T$ en $P$, of de ideale-gaswet.
  \item \textbf{Tweefasegebied} (verzadigd mengsel): $T$ en $P$ zijn \emph{niet} onafhankelijk, want
        $P = P_\mathrm{sat}(T)$. Je hebt een tweede eigenschap nodig, meestal de dampfractie $x$:
        \[
        h = h_f + x\,h_{fg},\qquad s = s_f + x\,s_{fg},\qquad x = \frac{h-h_f}{h_g-h_f}.
        \]
        ``Verzadigde vloeistof'' is $x=0$ en ``verzadigde damp'' is $x=1$.
  \item \textbf{Combinaties zoals ($P$, $s$) of ($P$, $h$)} zijn net zo goed twee onafhankelijke eigenschappen. Na een
        isentrope compressie ken je bijvoorbeeld $P_2$ en $s_2=s_1$. Daarmee zoek je $h_{2s}$ op.
\end{itemize}

\section{Ideaal gas, reëel gas en de compressibiliteitsfactor}
\begin{methode}[– Cengel H3 en H4]
Ideaal gas: $P v = R T$ met $R = R_u/M$ en $R_u = 8{,}314$\,kJ/(kmol$\cdot$K). Voor een ideaal gas hangt $h$
alleen van $T$ af. Met constante soortelijke warmte geldt $\Delta h = c_p\,\Delta T$ (tabel A-2), met
$k = c_p/c_v$. Voor een reëel gas schrijf je $P v = Z R T$. De compressibiliteitsfactor $Z$ volgt uit
gegeneraliseerde kaarten of uit de viriaalvergelijking $Z = 1 + B(T)\,P/(R_u T)$.
\end{methode}
Waar gebruiken we wat?
\begin{itemize}
  \item \OO\ en waterdamp bij 1,1 tot 1,7\,bar: de partiële drukken zijn zo laag dat ideaal gas uitstekend klopt
        ($Z$ wijkt minder dan 0,1\,\% af van 1).
  \item \HH\ bij 21 tot 71\,bar: $Z$ ligt tussen \VAR{n(RA.Z_in_K101,3)} en \VAR{n(RA.Z_V102_max,3)}. Dat is een
        afwijking van 1 tot 4\,\%. Versie B rekent toch met ideaal gas. Voor de buffervaten neemt versie B
        $Z$ mee via de viriaalcoëfficiënt. Versie A gebruikt de reële-gasvergelijking van CoolProp.
  \item Koudemiddelen: die zijn sterk niet-ideaal (tweefase, dicht bij het kritische punt). Daar is altijd een
        tabel of diagram nodig.
\end{itemize}

\section{De energiebalans voor stationaire stroming}
\begin{methode}[– Cengel H5, ``steady-flow energy balance'']
Voor een controlevolume in stationaire toestand (niets verandert in de tijd) geldt met één in- en één uitgang:
\[
\Qd - \Wd = \md\left[(h_2-h_1) + \tfrac{V_2^2-V_1^2}{2} + g(z_2-z_1)\right] \;\approx\; \md\,(h_2-h_1).
\]
De kinetische en potentiële energie verwaarlozen we. Bij gassnelheden onder 20\,m/s is
$V^2/2 < 0{,}2$\,kJ/kg, en dat valt weg tegen de $\Delta h$ van tientallen kJ/kg.
\end{methode}
De standaardapparaten uit Cengel H5 komen allemaal in het PFD voor:
\begin{center}
\begin{tabularx}{\textwidth}{lXl}
\toprule
apparaat (Cengel H5) & vereenvoudiging & in het PFD\\
\midrule
compressor & $\Qd\approx 0$ (adiabatisch), dus $\Wd_\mathrm{in} = \md(h_2-h_1)$ & K-101, K-102, K-201, K-202\\
pomp & $\Wd_\mathrm{in} = \md\,v\,\Delta P/\eta$ & P-101\\
warmtewisselaar & geen arbeid, dus $\Qd = \md(h_\mathrm{uit}-h_\mathrm{in})$ per stroom & E-101 t/m E-206\\
smoorventiel & $\Qd=\Wd=0$, dus $h_\mathrm{uit} = h_\mathrm{in}$ & TV-201, TV-202, PCV-101\\
mengkamer/splitsing & $\sum \md_\mathrm{in} h_\mathrm{in} = \sum \md_\mathrm{uit} h_\mathrm{uit}$ & TW-verdeling\\
\bottomrule
\end{tabularx}
\end{center}
\textbf{Tekenafspraak.} $\Qd$ is positief als er warmte \emph{naar} de processtroom gaat. Een koeler zoals E-101
heeft dus een \emph{negatieve} $\Qd$: de processtroom staat warmte af. Op de PFD staat ``E-101: \VAR{n(R.Q_E101,1)}\kW''.
Het minteken betekent dat er \VAR{n(-R.Q_E101,1)}\kW\ uit het gas wordt gehaald.

\section{Entropie, isentroop rendement en pompen}
\begin{methode}[– Cengel H7]
\textbf{Isentroop rendement van een compressor:}
$\ds \eta_C = \frac{\text{isentrope arbeid}}{\text{werkelijke arbeid}} = \frac{h_{2s}-h_1}{h_2-h_1}
\;\Rightarrow\; h_2 = h_1 + \frac{h_{2s}-h_1}{\eta_C}$.

\textbf{Isentroop proces van een ideaal gas met constante $c_p$:}
$\ds \frac{T_{2s}}{T_1} = \left(\frac{P_2}{P_1}\right)^{(k-1)/k}$.

\textbf{Entropie van een ideaal gas:}
$s_2-s_1 = c_p\ln\frac{T_2}{T_1} - R\ln\frac{P_2}{P_1}$.

\textbf{Pomp (onsamendrukbare vloeistof):}
$w_s = v\,(P_2-P_1)$ en $\Wd = \md\,v\,\Delta P/\eta_P$.
\end{methode}
\begin{waarom}[isentroop als referentie]
Een adiabatische compressie zonder wrijving is omkeerbaar, en dan is $s$ constant. Dat is de beste compressie die
mogelijk is: de minste arbeid. Een echte machine heeft verliezen (wrijving, lekkage, turbulentie). Die verliezen
worden extra warmte in het gas, dus $h_2 > h_{2s}$ en $T_2 > T_{2s}$. Het rendement $\eta$ is een ontwerpgetal
van de leverancier, typisch 0,65 tot 0,85. Daarom rekenen we altijd \emph{eerst} de isentrope eindtoestand uit en
corrigeren dan met $\eta$.
\end{waarom}

\section{Gasmengsels en gas-dampmengsels}
\begin{methode}[– Cengel H13 en H14]
Molfractie $y_i = N_i/N_\mathrm{tot}$ en $N_i = m_i/M_i$. \textbf{Wet van Dalton:} de partiële druk is
$P_i = y_i\,P$. Elke component gedraagt zich alsof hij alleen bij $T$ en $P_i$ in het volume zit. Voor een
mengsel geldt $h_\mathrm{mix} = \sum m_i h_i/m$ en $c_{p,\mathrm{mix}} = \sum m_i c_{p,i}/m$. Voor de
entropie van elke component reken je bij zijn \emph{partiële} druk.

\textbf{Gas-dampmengsel} (H14, bij lucht en waterdamp, hier bij \OO\ en waterdamp): de damp condenseert zodra
de partiële dampdruk de verzadigingsdruk bereikt. Het \emph{dauwpunt} is $T_\mathrm{dauw} = T_\mathrm{sat}(P_v)$.
Na afkoelen tot $T$ is het gas verzadigd en geldt $P_v = P_\mathrm{sat}(T)$. Het overschot aan water condenseert.
De vochtverhouding is $\omega = m_v/m_\mathrm{gas}$. Voor lucht is dat
$0{,}622\,P_v/(P-P_v)$; voor \OO\ is de factor $M_{\HHO}/M_{\OO} = 0{,}563$.
\end{methode}

\section{Warmtepompen: COP, Carnot en $\eta_\mathrm{II}$}
\begin{methode}[– Cengel H6, H8 en H11]
Een warmtepomp neemt warmte $Q_L$ op bij lage temperatuur $T_L$. Met arbeid $W$ geeft hij $Q_H = Q_L + W$ af
bij hoge temperatuur $T_H$. Dan geldt:
\[
\mathrm{COP}_\mathrm{HP} = \frac{Q_H}{W},\qquad
\mathrm{COP}_\mathrm{HP,Carnot} = \frac{T_H}{T_H-T_L}\ (T\text{ in K}),\qquad
\eta_\mathrm{II} = \frac{\mathrm{COP}}{\mathrm{COP}_\mathrm{Carnot}}.
\]
Bij de dampcompressiecyclus (H11) lees je per kg koudemiddel af:
$q_L = h_1-h_4$, $w = h_2-h_1$ en $q_H = h_2-h_3$. Dan is $\mathrm{COP} = (h_2-h_3)/(h_2-h_1)$.
\end{methode}
\begin{waarom}[een warmtepomp en niet direct verwarmen]
Volgens de tweede hoofdwet (Clausius, Cengel H6) stroomt warmte vanzelf alleen van warm naar koud. Het
koelwatercircuit TW is hooguit \VAR{n(R.T_tw_a,1)}\degC. Het kan het netwater daarom nooit tot 70\degC\ verwarmen.
Daar is arbeid voor nodig, en dat is de warmtepomp. Het netwater komt wel terug met 40\degC, en dat is
\emph{kouder} dan TW. Dat stuk (40 naar 45\degC) gaat zonder arbeid via E-204.
\end{waarom}

\section{Verbrandingswaarde en elektrolyse}
\begin{methode}[– Cengel H15, tabel A-27]
De \emph{hogere verbrandingswaarde} (HHV) van waterstof is 141\,800\,kJ/kg; daarbij is het product
\HHO\ vloeibaar. De lagere verbrandingswaarde (LHV) is 120\,000\,kJ/kg; daarbij is het product \HHO\ damp.
Elektrolyse is de omgekeerde reactie, $\HHO(\ell) \rightarrow \HH + \tfrac12\OO$. Die vraagt dus minstens de
reactie-enthalpie, en dat is de HHV.
\end{methode}

%% if A
\section{Wat CoolProp doet, in Cengel-taal}
\label{s:coolprop}
CoolProp is een open-source bibliotheek met \emph{toestandsvergelijkingen} van hoge nauwkeurigheid: per stof een
formule voor de vrije Helmholtz-energie $a(T,\rho)$, waaruit alle eigenschappen volgen. Het zijn dezelfde
referentievergelijkingen waarmee de tabellen in Cengel zijn gemaakt. Cengel heeft zijn tabellen met het programma
EES uit deze vergelijkingen laten genereren. Voor water is dat IAPWS-95.
\begin{center}
\begin{tabularx}{\textwidth}{lX}
\toprule
stof & toestandsvergelijking in CoolProp\\
\midrule
water & IAPWS-95 (zelfde basis als Cengel tabel A-4 t/m A-7)\\
waterstof & Leachman e.a. (2009), \emph{normal hydrogen}\\
zuurstof & Schmidt \& Wagner (1985)\\
R1234ze(E) & Thol \& Lemmon (2016)\\
R1233zd(E) & Mondéjar e.a. (2015)\\
\bottomrule
\end{tabularx}
\end{center}
\textbf{Een CoolProp-aanroep is een tabelopzoeking.} De aanroep
\[
\texttt{PropsSI("H", "T", 303.15, "P", 107078, "Oxygen")}
\]
betekent in Cengel-taal: ``zoek de enthalpie $h$ op van zuurstof bij $T = 30\degC$ en $P = 107{,}1\kPa$''. Dat is het
toestandspostulaat met de twee onafhankelijke eigenschappen $T$ en $P$. CoolProp rekent in SI-basiseenheden
(K, Pa, J/kg). Het model deelt daarna door 1000 voor kJ/kg en vermenigvuldigt met 1000 voor kPa.
Veelgebruikte invoerparen in dit verslag:
\begin{center}
\begin{tabular}{lll}
\toprule
invoer & betekenis & Cengel-equivalent\\
\midrule
\texttt{"T", T, "Q", 0} & verzadigde vloeistof bij $T$ & tabel A-4: $h_f$, $P_\mathrm{sat}$\\
\texttt{"T", T, "Q", 1} & verzadigde damp bij $T$ & tabel A-4: $h_g$\\
\texttt{"P", P, "Q", 1} & verzadigde damp bij $P$ & tabel A-5: $T_\mathrm{sat}$\\
\texttt{"T", T, "P", P} & eenfasegebied & tabel A-6/A-7, of ideaal gas\\
\texttt{"P", P, "S", s} & na isentrope compressie & tabel A-6 interpoleren op $s$\\
\texttt{"P", P, "H", h} & na compressie met $\eta$ of smoren & tabel A-6 interpoleren op $h$\\
\bottomrule
\end{tabular}
\end{center}
\begin{letop}[nulpunt van $h$ en $s$]
Absolute waarden van $h$ en $s$ hangen af van een gekozen nulpunt. CoolProp gebruikt voor koudemiddelen het
IIR-nulpunt ($h=200$\,kJ/kg en $s=1$\,kJ/(kg$\cdot$K) voor verzadigde vloeistof bij 0\degC). Cengel gebruikt
voor R-134a $h=s=0$ voor verzadigde vloeistof bij $-40\degC$. Voor water gebruiken beide het tripelpunt. Je
mag dus \emph{nooit} een $h$ uit CoolProp mengen met een $h$ uit een tabel met een ander nulpunt. Verschillen
$\Delta h$ zijn wel altijd gelijk, en alle balansen werken met verschillen.
\end{letop}
%% else
\section{Wat ``alleen Cengel'' hier betekent}
In deze versie komen alle stofeigenschappen uit de tabellen van Cengel:
\begin{itemize}
  \item tabel A-1: molaire massa's en gasconstanten ($R_{\HH} = 4{,}124$, $R_u = 8{,}314$);
  \item tabel A-2: $c_p$ en $k$ van \OO, \HH\ en waterdamp bij 300\,K (constante soortelijke warmte);
  \item tabel A-4: verzadigd water (temperatuurtabel): $P_\mathrm{sat}$, $v_f$, $h_f$, $h_g$;
  \item tabel A-27: HHV van waterstof.
\end{itemize}
Voor de koudemiddelen bestaan geen Cengel-tabellen. Die waarden lezen we af uit het log $p$–$h$-diagram
(figuren in hoofdstuk~\ref{h:wp}). De afgelezen waarden zijn afgerond op 0,5\,kJ/kg. Dat is de nauwkeurigheid
waarmee je op een groot afgedrukt diagram kunt aflezen.
\begin{letop}[nulpunt van $h$]
De diagrammen gebruiken het IIR-nulpunt ($h=200$\,kJ/kg voor verzadigde vloeistof bij 0\degC). De
absolute getallen zijn daardoor niet te vergelijken met de R-134a-tabellen in Cengel, die $h=0$ bij $-40\degC$
nemen. Alleen verschillen $\Delta h$ tellen, en die zijn onafhankelijk van het nulpunt.
\end{letop}
%% endif
